\documentclass[a4paper,11pt,twoside,openright]{book}

\usepackage[spanish]{babel}
\usepackage[utf8]{inputenc}

\usepackage{eurosym}
\usepackage{fancybox}
\usepackage{multicol}
\usepackage{amsmath}
\usepackage{subcaption}
\usepackage[nottoc,numbib]{tocbibind}
\usepackage[pdftex,usenames,dvipsnames]{color}

% compilar con pdflatex.
\usepackage[pdftex]{graphicx}


\usepackage{setspace}
\onehalfspacing
\usepackage{geometry}
\geometry{
  inner=3.5cm,
  outer=2.5cm,
  bottom=3.5cm,
  top=3.5cm}

\usepackage{fancyhdr}
\pagestyle{fancy}
\fancyhf{}
\fancyhf[HR]{\thepage}
\fancyhf[HL]{\nouppercase\rightmark}

\usepackage{booktabs}
\usepackage{hyperref}
\hypersetup{
    colorlinks,
    linkcolor={red!50!black},
    citecolor={blue!50!black},
    urlcolor={blue!80!black}
}
\usepackage{rotating}
\usepackage{multicol}
\usepackage{multirow}
\usepackage{pgfgantt}
\usepackage{enumitem}
\usepackage{color}
\usepackage{xcolor}
\usepackage{caption}
\DeclareCaptionFont{white}{\color{white}}
\DeclareCaptionFormat{listing}{\colorbox{gray}{\parbox{\textwidth}{#1#2#3}}}
\captionsetup[lstlisting]{format=listing,labelfont=white,textfont=white}

 \usepackage{listings}
 \usepackage{courier}
 \lstset{
         basicstyle=\footnotesize\ttfamily, % Standardschrift
         numberstyle=\tiny,          % Stil der Zeilennummern
         numbersep=5pt,              % Abstand der Nummern zum Text
         tabsize=2,                  % Groesse von Tabs
         extendedchars=true,         %
         breaklines=true,            % Zeilen werden Umgebrochen
         keywordstyle=\textbf,
         frame=b,
         stringstyle=\textit, % Farbe der String
         showspaces=false,           % Leerzeichen anzeigen ?
         showtabs=false,             % Tabs anzeigen ?
         xleftmargin=17pt,
         framexleftmargin=17pt,
         framexrightmargin=5pt,
         framexbottommargin=4pt,
         %backgroundcolor=\color{lightgray},
         showstringspaces=false      % Leerzeichen in Strings anzeigen ?
 }

 \lstset{literate=
  {á}{{\'a}}1 {é}{{\'e}}1 {í}{{\'i}}1 {ó}{{\'o}}1 {ú}{{\'u}}1
  {Á}{{\'A}}1 {É}{{\'E}}1 {Í}{{\'I}}1 {Ó}{{\'O}}1 {Ú}{{\'U}}1
  {à}{{\`a}}1 {è}{{\`e}}1 {ì}{{\`i}}1 {ò}{{\`o}}1 {ù}{{\`u}}1
  {À}{{\`A}}1 {È}{{\'E}}1 {Ì}{{\`I}}1 {Ò}{{\`O}}1 {Ù}{{\`U}}1
  {ä}{{\"a}}1 {ë}{{\"e}}1 {ï}{{\"i}}1 {ö}{{\"o}}1 {ü}{{\"u}}1
  {Ä}{{\"A}}1 {Ë}{{\"E}}1 {Ï}{{\"I}}1 {Ö}{{\"O}}1 {Ü}{{\"U}}1
  {â}{{\^a}}1 {ê}{{\^e}}1 {î}{{\^i}}1 {ô}{{\^o}}1 {û}{{\^u}}1
  {Â}{{\^A}}1 {Ê}{{\^E}}1 {Î}{{\^I}}1 {Ô}{{\^O}}1 {Û}{{\^U}}1
  {œ}{{\oe}}1 {Œ}{{\OE}}1 {æ}{{\ae}}1 {Æ}{{\AE}}1 {ß}{{\ss}}1
  {ç}{{\c c}}1 {Ç}{{\c C}}1 {ø}{{\o}}1 {å}{{\r a}}1 {Å}{{\r A}}1
  {€}{{\EUR}}1 {£}{{\pounds}}1
 }

\parskip=6pt

\parindent=10pt

\setcounter{secnumdepth}{3}

\title{Desarrollo de un videojuego roguelike para invidentes aplicando técnicas de Procesamiento del Lenguaje Natural.}
\author{Darío Penas Sabín}
\date{\today}


\renewcommand\lstlistingname{Fragmento de Código}
\renewcommand\lstlistlistingname{Fragmentos de Código}

\newenvironment{bottompar}{\par\vspace*{\fill}}{\clearpage}

\begin{document}
        %
% Portada.
%

% Nota: Sería más cómodo emplear el comando \maketitle que genera una portada de forma automática, pero
% no incluye toda la información que es necesario incluir en la memoria de un proyecto de fin de carrera
% de la Facultad de Informática de A Coruña.
%

\begin{titlepage}

	\begin{center}

		% Logotipo de la universidad.
		\includegraphics[width=6cm]{./eps/logo_udc.eps}
		\vspace{2cm}

		{\Large{\textbf{Facultade de Informática da Universidade de A Coruña}}}
		\\
		{\it \large{\textbf{Computación}}}
		\vspace{1cm}

		% Indicamos el nombre de la titulación oficial que hemos cursado con tanto esfuerzo.
		{\large PROYECTO DE FIN DE CARRERA\\Ingeniería Informática}
		\vspace{1cm}

		% Título
		\textbf{\Large Desarrollo de un videojuego roguelike para invidentes aplicando técnicas de Procesamiento del Lenguaje Natural.}
		\vspace{7cm}
	\end{center}

	\begin{flushright}
		\begin{tabular}{ll}
			\large{\textbf{Alumno:}}	&
			\large{Darío Penas Sabín} \\

			% Nombre del director/tutor del proyecto.
			\large{\textbf{Director:}}	&
			\large{Jesús Vilares Ferro} \\

			% Nombre del director/tutor del proyecto.
			\large{\textbf{Director:}}	&
			\large{Carlos Gómez Rodríguez} \\

			% Fecha.
			\large{\textbf{Fecha:}}	&
			\large{CHANGE: 15 de junio de 2016} \\
		\end{tabular}
	\end{flushright}

\end{titlepage}


        \frontmatter

        \thispagestyle{empty}

        \section*{Resumen:}

La industria del entretenimiento digital ha crecido inmensamente en los últimos años, llegando a alcanzar números de ventas jamás vistos anteriormente.
Parte de la razón de este crecimiento viene dada por una mejora radical en el aspecto visual, necesaria para que el jugador se sienta inmerso en la aventura que se le está planteando.
Estas mejoras, sin embargo, dejan de lado a muchos jugadores que, por diferentes motivos, no son capaces de apreciar el contenido visual que se les ofrece o tienen problemas para ello, haciendo imposible que disfruten del contenido ofrecido.

Este proyecto consiste en la creación de un videojuego de género \textit{roguelike} para invidentes que, desde un principio, parte de la idea de generar contenido específicamente diseñado para que pueda ser jugado por todo el mundo, haciendo énfasis en ofrecer al jugador una diversa cantidad de frases que describan lo que está sucediendo en su alrededor y que serán generadas automáticamente en base a una serie de gramáticas y diccionarios, empleando para ello técnicas de Procesamiento del Lenguaje Natural.

        \thispagestyle{empty}

\end{document}

